% =============================================================================
% First-year progression report
% Topological Feature Learning for Parameter Estimation in Neyman-Scott
% Processes -- Flavio Gualtieri
%
% Text in blue (\new{...}) is NOT lifted verbatim from the preprint.
% Text in red (\rev{...}) is new or altered in this revision.
% =============================================================================

\documentclass[11pt,a4paper]{article}

\usepackage[margin=1in]{geometry}
\usepackage{mathpazo}
\usepackage{microtype}

\usepackage{amsmath}
\usepackage{amssymb}
\usepackage{amsthm}

\usepackage{booktabs}
\usepackage{enumitem}
\usepackage[font=small,labelfont=bf,skip=4pt]{caption}
\setlength{\tabcolsep}{5pt}
\renewcommand{\arraystretch}{1.0}
\setlength{\abovedisplayskip}{5pt}
\setlength{\belowdisplayskip}{5pt}
\setlength{\abovedisplayshortskip}{3pt}
\setlength{\belowdisplayshortskip}{3pt}

\usepackage{natbib}
\usepackage{xcolor}
\usepackage[hidelinks]{hyperref}
\raggedbottom
\renewcommand{\textfraction}{0.05}
\renewcommand{\topfraction}{0.95}
\renewcommand{\bottomfraction}{0.95}
\renewcommand{\floatpagefraction}{0.85}

% ----------------------------------------------------------------------------
% Macros
% ----------------------------------------------------------------------------
\newcommand{\R}{\mathbb{R}}
\newcommand{\N}{\mathbb{N}}
\newcommand{\E}{\mathbb{E}}
\newcommand{\pc}{\mathbf{x}}
\newcommand{\simpcom}{\Sigma}
\newcommand{\filt}{\mathcal{F}}
\newcommand{\kpar}{\kappa}
\newcommand{\mo}{\mu}

% ----------------------------------------------------------------------------
\definecolor{newtextcolor}{rgb}{0.00,0.35,0.70}
\newcommand{\new}[1]{\textcolor{newtextcolor}{#1}}
\newcommand{\rev}[1]{\textcolor{red}{#1}}

% ----------------------------------------------------------------------------
\newtheoremstyle{definitioncolon}%
  {4pt}{4pt}{\normalfont}{}{\bfseries}{}{.5em}%
  {\thmname{#1}\thmnumber{ #2}\thmnote{: #3}}
\theoremstyle{definitioncolon}
\newtheorem{definition}{Definition}[section]

\newtheorem*{note}{Note}
\newtheorem*{example}{Example}

% ----------------------------------------------------------------------------
\title{Topological Feature Learning for Parameter Estimation\\
       in Neyman-Scott Processes\\[4pt]
       \large \new{First-year progression report}}
\author{Flavio Gualtieri}
\date{\today}

\begin{document}
\maketitle



% #############################################################################
\section{Introduction}
\label{sec:intro}
% #############################################################################

The study of spatial point processes finds applications in astronomy, cell
biology, epidemiology, etc \citep{moller2004statistical}. Many natural phenomena
-- such as galaxy clustering -- are modelled using SPPs REF, so one can hope to
leverage an understanding of these processes to obtain novel insights. The
parameters of spatial point processes are of particular interest because
they are interpretable scientific quantities.
For example, when Neyman-Scott processes are used to model galaxy clustering the parent intensity $\kappa$ corresponds to the number density of
haloes hosting galaxies, the offspring intensity $\mu$ is the mean occupancy of
a halo, and the cluster scale $\sigma$ is the spatial extent of the galaxy
distribution within a halo; these are exactly the quantities a halo occupancy
model is built to constrain REF. In forest ecology,
\citet{waagepetersen2009twostep} one can fit Thomas processes to three tropical tree
species and read the cluster scale directly as a summary of how tightly seeds
are dispersed, finding that species dispersed by exploding capsules, by wind,
and by birds and mammals separate along that parameter.

A number of parameter estimation methods have been proposed,
typically relying on closed-form expressions for summary statistics, such as the $K$ and $L$
functions (see Section~\ref{ssec:spp}). Minimum contrast estimation fits
parameters by numerically
minimizing a discrepancy between the empirical and theoretical curve, most
commonly for $K$ or $g$ \citep{diggle2003spatial}. Palm likelihood instead
builds a composite likelihood from the process's Palm intensity, and has been
developed specifically for Neyman-Scott processes \citep{tanaka2008parameter}.
Fully Bayesian approaches treat the unobserved parent locations as latent
variables and sample the posterior over $\theta$ by MCMC
\citep{moller2004statistical}.

These methods have three limitations we hope to address with our proposed model. Firstly, they rely on the existence of closed-form expressions of
summary statistics which may not exist for every process.
This causes two main issues: (i) The expression may not exist: the
LGCP-Strauss process of \citet{vihrs2022neural} does not have a tractable density, and (ii) The expression may exist but fail
the regularity conditions the estimator's theory needs:
\citet{waagepetersen2009twostep} require $K_\psi$ to be twice continuously
differentiable in $\psi$, which explicitly excludes the Mat\'ern cluster
process.

Secondly, classical parameter estimation methods suffer from
hyperparameter sensitivity, specifically the minimum contrast method
depends on $(r_{max}, q, p)$. Naturally, ParamNet also has hyperparameters, but we claim that they are more easily handled and optimized. Minimum contrast's $(r_{max},q,p)$ are chosen per dataset, by rule of thumb, with
no held-out quantity that tells the practitioner whether the choice was optimal.
\citet{waagepetersen2009twostep} report that estimates of $\kappa$ and $\omega$
for a Thomas process are ``quite sensitive'' to the choice of $r_{max}$ and $q$,
and attribute this to a strong negative correlation between $\hat\kappa$ and
$\hat\omega$ -- the product $\hat\kappa\hat\omega$ is stable while the
individual estimates are not. ParamNet's hyperparameters are fixed once on a
training split, selected against validation loss, then frozen and applied
unchanged to every future cloud. The choice is therefore made once, with a
measurable criterion, rather than repeated per dataset with none.

Finally, classical methods become unstable as the process approaches complete
spatial randomness (near-Poisson regime). Let $\lambda=\kappa\mu$ be the expected number of points per unit area, and $\nu = \sigma\sqrt{\kappa}$ the ratio of cluster radius to mean spacing between parents (both are defined properly in
Section~\ref{ssec:simulator}). When $\nu \ll 1$ the clusters are well separated
and you can practically count them, so $\kappa$ is easy. As $\nu$ grows past
$1$, neighbouring clusters overlap and the pattern looks Poisson: many
$(\kappa,\mu)$ pairs with the same product $\lambda$ produce nearly the same
observable pattern, so $\kappa$ and $\mu$ are only weakly separately
identifiable from a single realization. The theoretical $K$-function makes this
explicit -- for a Thomas process $K(t) = \pi t^2 + (1-e^{-t^2/4\sigma^2})/\kappa$,
so as $\sigma$ grows at fixed $t$ the second term vanishes and $K$ tends to the
Poisson $K$ regardless of $\kappa$ \citep{waagepetersen2009twostep}.
Minimum contrast is fitted by numerical optimization of
a discrepancy that is close to flat along this trade-off, so the estimates it
returns in this regime are sensitive to the starting point and to the choice of
$r_{max}$.

\citet{waagepetersen2009twostep} note this failure in their simulations:
in their weakly clustered case the estimated $K$ could not be distinguished from
the Poisson $K$, and ``rather extreme'' clustering estimates resulted.

We deliberately retain this regime in our design distribution (see
Section~\ref{ssec:simulator}) rather than restricting to well-separated
clusters.

These limitations motivate the use of persistence-based features which provide a
multi-scale, process-agnostic descriptor of the point cloud that does not require
a closed-form or a choice of summary statistic.

% -----------------------------------------------------------------------------
\subsection{Problem setup}
\label{ssec:setup}
% -----------------------------------------------------------------------------

Let $X$ be a Neyman-Scott process on $\R^d$ with unknown parameter vector
$\theta \in \R^{n_{params}}$. For example, the homogeneous Thomas process has
$\theta = (\kappa, \mu, \sigma)$ corresponding to parent intensity, offspring intensity and 
cluster scale, respectively. We observe a single realization of $X$ on a
bounded window $W \subset \R^d$, giving a point cloud
$\pc = \{x_1, \dots, x_n\} \subset \R^d$ with $n = |X \cap W|$. The objective is to
produce an estimate $\hat\theta = \hat\theta(\pc)$ of $\theta$ from $\pc$ alone.
We stress \emph{single realization}: $\theta$ is not recovered from repeated
draws of $X$. This is the setting a practitioner
faces: one observed pattern and one chance to estimate the generating parameters.
It is also what makes the problem hard in a way no estimator can fix: a
fixed $\theta$ already induces a distribution over point clouds, and it is this
irreducible
spread -- not any weakness of a particular estimator -- that sets the error floor
characterized below.

The model is trained on a set of synthetic clouds
$\{\pc_i, \theta_i\}_{i=1}^{N_{train}}$ by backpropagating the RMS loss
    \[
    \mathcal{L}(\hat\theta_i, \theta_i) = \sqrt{\sum_{j=1}^{n_{params}}(\hat\theta^j_i - \theta^j_i)^2}.
    \]
The scale of each
parameter varies; therefore -- if not normalized -- the
large-scale parameters will
dominate this loss. To address this, we train in log-normalized parameter space,
mapping $\theta_i \mapsto (\log(\theta_i) - \mu_{\log\theta})/\sigma_{\log\theta}$,
where -- abusing notation -- $\mu_{\log\theta}, \sigma_{\log\theta}$ denote the
mean and standard deviation of the logged parameters over the training set.

% -----------------------------------------------------------------------------
% Commented out for now (2026-08-21): the report stands without the
% irreducible error floor. Re-enable by deleting the \iffalse/\fi pair
% below once L_floor actually has a value (see the label-indexing NaN bug
% noted in the -- now also commented-out -- Outstanding-experiments bullet
% near the end of the document).
% -----------------------------------------------------------------------------
\iffalse
\paragraph{The irreducible error floor.}
We simulate one realization per parameter vector, so the training set never
shows the model two clouds with the same $\theta$. The quantity that acts as an error floor is the \emph{Bayes risk} of the estimation problem under
the design distribution $\pi(\theta)$,
\[
  \mathcal{L}_{\text{floor}}
  \;=\;
  \E_{\theta \sim \pi}\,\E_{\pc \mid \theta}
  \big[\, \|\,\E[\theta \mid \pc] - \theta\,\|^2 \,\big]
  \;=\;
  \E_{\pc}\big[\operatorname{tr}\operatorname{Var}(\theta \mid \pc)\big],
\]
in the same log-normalized coordinates as the reported losses. This is the mean
squared error of the best possible estimator, and no estimator seeing a single
realization can beat it on average over the same design distribution.

Two routes are available to us, both
simulation-only: (i) fix a grid of $\theta$, draw $R$ replicate realizations at
each, and report the replicate dispersion of \emph{any} fixed per-realization
estimator as a benchmark -- cheap, but estimator-dependent; (ii) fit a
simulation-based posterior (e.g.\ a conditional density estimator, or an
ensemble trained with a proper scoring rule) on the same design distribution and
average its posterior variance -- more faithful to the definition above. We
report route (i) first because it is a one-line change to existing code, and
flag route (ii) as the correct target.
\rev{Note that neither route needs a closed form; both need only a simulator.}
\fi

% #############################################################################
\section{Background}
\label{sec:background}
% #############################################################################

\citet{vihrs2022neural} first proposed the use of deep learning for parameter
estimation. They implement
an encoder/inference architecture similar to PointNet, but use the $L$-function as the learned
feature.
\citet{yip2025learning} employ a
similar architecture as well, but use
persistence images as the learned feature,mapping $H_0$--$H_2$ images of an
$\alpha\mathrm{DTM}\ell$ filtration of a halo catalogue to $(\Omega_m,\sigma_8)$.
Their model is designed to estimate cosmological parameters rather than the
parameters of the underlying point process.
Furthermore, they do not systematically compare
various topological summaries, nor do they address the calibration of the persistence images: they grid-search the kernel bandwidth and raster
resolution, but the birth and persistence axis ranges -- which fix what a pixel
means -- are not treated as a design choice at all.

The vectorization of persistence diagrams is an active area of research arising from the issue
that persistence diagrams -- not being a Hilbert space -- cannot be directly used as inputs for
machine learning models. To mitigate this, \citet{adams2017persistence} propose persistence
images, which vectorise persistence
diagrams by applying a Gaussian kernel to each point in the diagram and integrating over a grid.
\citet{umeda2017time} proposes Betti curves, which vectorize persistence diagrams by counting
the number of topological features
that are alive at each scale.
\citet{bubenik2015landscapes} proposes persistence landscapes and
\citet{chazal2014silhouette} adapts this idea into persistence silhouettes, both stacking ``tent'' functions over the diagram points; \citet{reininghaus2015kernel} avoid vectorization altogether with a kernel on diagrams.

The remainder of this section covers the basic notions in spatial point process and
persistent homology that the later sections assume.

% -----------------------------------------------------------------------------
\subsection{Spatial point processes}
\label{ssec:spp}
% -----------------------------------------------------------------------------

    \begin{definition}[Neyman-Scott process]
    \label{def:ns}
        A $d$-dimensional SPP $X$ on $W \subseteq \R^d$ is equipped
        with an intensity function $\rho(u): \R^d \rightarrow \R^+$ such that
        $\E[N(W)] = \int_W \rho(u) du$,
        where $N(W) = \mid W \cap X \mid$ denotes the number of points of $X$
        in $W$. Let $P$ be a homogeneous Poisson \emph{parent} process with
        intensity $\kpar$. For a Neyman-Scott process with Poisson-distributed
        offspring we have
            $\rho_{NS}(u) = \mo \sum_{p \in P} k(u - p)$,
        where $k$ is a kernel that determines how the \emph{offspring} points
        are distributed around the parents.
    \end{definition}

    \begin{note}
        Neyman-Scott processes always have the parent and offspring intensities
        as parameters. Additional parameters are introduced by the choice of
        kernel. A homogeneous Thomas process is a Neyman-Scott process with an
        isotropic Gaussian kernel $\mathcal{N}(0, \sigma^2I)$, and has
        parameters $\theta_{Thomas} = (\kappa, \mu, \sigma)$ corresponding to
        the parent intensity, offspring intensity, and cluster scale,
        respectively. For stationary and isotropic $X$, the pair correlation
        function $g(r) = \rho^{(2)}(r)/\rho^2$ and the $K$-function
        $K(r) = \int_{b_r(0)} g(\|h\|) dh$ are the two second-order summaries
        the classical estimators above are fitted against: $\rho(u)K(r)$ is the
        expected number of points we expect to
        see within a distance $r$ of a typical point in $X$, and
        $g$ can suggest whether SPPs are attractive ($g > 1$), or repulsive
        ($g < 1$).
        For the Thomas process
        $g(h) = 1 + e^{-\|h\|^2/4\sigma^2}/(4\pi\sigma^2\kappa)$ and
        $K(t) = \pi t^2 + (1-e^{-t^2/4\sigma^2})/\kappa$
        \citep{waagepetersen2009twostep}. Neither depends on $\mu$, which enters
        only through $\rho=\kappa\mu$; therefore any $K$- or $g$-based estimator cannot discern ...
    \end{note}

% -----------------------------------------------------------------------------
\subsection{Persistent homology}
\label{ssec:ph}
% -----------------------------------------------------------------------------

Radially averaged statistics -- such as the $g$ and $K$ functions -- collapse
a point cloud to a real-valued curve indexed by scale $r$, discarding
multi-scale connectivity information and
void structure, i.e. \emph{how} the density is distributed across the space.
Persistent homology retains this information by tracking a family of simplicial
complexes on the point cloud as the scale parameter grows. The $H_0$ persistence
summary records
the order and scales at which clusters merge; the $H_1$ counterpart marks voids and gaps
in the pattern as they persist across scales.
For each $h \in \N$, the $h$-dimensional persistence summary requires no closed form expression
and no assumption about which process family generated it.

Let $\pc = \{x_1, ..., x_n\} \subset \R^d$ be the point cloud of $X \cap W$, and
let $\Sigma(\pc)$ denote the \emph{full
simplex} on $\pc$, i.e. the collection of all non-empty subsets of $\pc$.

    \begin{definition}[Filtration]
    \label{def:filtfunc}
        A filtration function on $\pc$ is a monotone map
        $f: \Sigma(\pc) \rightarrow \R^+$ such that
        $f(\tau) \leq f(\sigma)$ for all $\tau \subseteq \sigma$.
        Given $r \geq 0$ we construct
        $\filt_r = \{\sigma \in \Sigma(\pc) : f(\sigma) \leq r\}$;
        monotonicity of $f$ ensures each $\filt_r$ is a simplicial complex
        and that $\filt_s \subseteq \filt_t$ for $0 \leq s \leq t$.
    \end{definition}

    \begin{example}
    \label{ex:dtm}
        For $u \in \R^d$, let $N_k(u)$ denote its $k$ nearest neighbours in
        $\pc$. The $\mathrm{DTM}_k$ filtration function is
        $f_{DTM}(\sigma; k) = \max_{u \in \sigma}\{(\frac{1}{k}
        \sum_{v \in N_k(u)} d(u, v)^2)^{\frac{1}{2}}\}$,
        extended to higher simplices by a max-over-edges rule
        \citep{anai2020dtm}. It is monotone, and is hence a valid
        filtration function \citep[Prop.\ 3.5]{anai2020dtm}.
    \end{example}

    \begin{definition}[Persistence diagram]
    \label{def:pd}
        A $h$-dimensional homology
        class is \emph{born} at the smallest scale $b \geq 0$ at which it
        appears in $H_h(\simpcom_b)$, and \emph{dies} at the scale $d > b$ at
        which it merges into an older class or becomes a boundary; $d - b$ is
        its \emph{lifetime}. A persistence diagram $D_h = \{(b_i, d_i)\}_i$ is
        the multiset of dimension-$h$ persistence pairs.
    \end{definition}

    \begin{definition}[Persistence image]
    \label{def:per_image}
        The persistence image of $D_h$ is constructed by mapping to
        birth-persistence coordinates $(b,d) \mapsto (b, d-b)$, weighting each
        point by $w(d-b)$, forming the persistence surface
        $\rho(z) = \sum_i w(d_i - b_i)\,\phi(z; (b_i, d_i - b_i))$ -- a
        weighted sum of isotropic
        Gaussian kernels with bandwidth $\sigma_{im}$ -- and integrating
        $\rho$ over a $64 \times 64$ pixel grid.
    \end{definition}


% #############################################################################
\section{ParamNet}
\label{sec:progress}
% #############################################################################

The work completed so far consists of the simulator and its validation, the
ParamNet model and its various architectural ablations, comparison with alternative
vectorizations, and benchmarking against existing baselines trained on the same
data.

% -----------------------------------------------------------------------------
\subsection{Synthetic Cloud Generation and experimental design}
\label{ssec:simulator}
% -----------------------------------------------------------------------------

Training data consists of pairs $(\pc_i, \theta_i)$ in which $\theta_i = (\theta_i^1, \dots, \theta_i^{n_{params}})$ is drawn
from a design distribution over parameters space and each $\pc_i$ is a
single realization of the corresponding process.
The naive approach of sampling $(\theta_i^1, \dots, \theta_i^{n_{params}})$ uniformly and
independently from per-parameter intervals is inadequate: point processes can
typically be separated into qualitative regimes of interest, which occupy curved
subsets of the parameter grid.
Therefore uniform sampling would allocate most of
its mass to configurations that are either near-Poisson or degenerate.
We therefore sample in reparameterized coordinates that carry the same meaning
for every family considered:
the expected offspring intensity $\lambda = \kappa\mu$, the parent intensity
$\kappa$, and the dimensionless overlap index $\nu = r\sqrt{\kappa}$,
where $r$ denotes the cluster scale of the family in question: $r = \sigma$ for
Thomas and $r = \sqrt{\sigma_1^2 + \sigma_2^2}$
for nested Thomas. The overlap index $\nu$ measures cluster radius against mean
inter-parent spacing, so $\nu \ll 1$ gives well-separated clusters and
$\nu \gtrsim 1$ gives clusters that merge. Nested Thomas carries an additional
coordinate, the scale ratio $\rho = \sigma_2/\sigma_1 \in (0,1)$.
This approach is general to the Neyman-Scott family: every realization
has parameters $\kappa$ and $\mu$ by Definition~\ref{def:ns}, and every kernel we
consider is a scale family $k(\cdot)=r^{-d}k_0(\cdot/r)$, so it supplies exactly
one length. Taking $r$ to be the root-mean-square displacement makes the
coordinates comparable across realizations.

We draw $(\log\kappa, \log\lambda, \log\nu)$ uniformly over a grid and invert to
recover $\mu = \lambda/\kappa$ and $r = \nu/\sqrt{\kappa}$ for the training labels;
log-uniform sampling
is used because the parameters span several orders of magnitude.
Three constraints are applied to the design grid. 
\begin{enumerate}
    \item \emph{Point budget}: the cost
    of computing persistence diagrams
    grows rapidly with the number of points, so we impose
    $\E[N] = \lambda \in [150, 800]$, fixing the range of
    $\log\lambda$.
    \item \emph{Clusters per window}: we require $\kappa \geq 15$
    so that a single realization is informative about $\kappa$, and
    $\kappa \leq 120$ so that individual clusters remain
    resolvable.
    \item \emph{Overlap}: we impose $\nu \in [0.10, 0.90]$, the upper
    limit deliberately admitting the near-Poisson regime so that the model is
    still exposed to configurations in which the parameters are weakly
    identifiable.
    Parameter draws violating these constraint are rejected and redrawn before any
    cloud is generated; this filter acts on the design and not on realizations,
    therefore it does not distort the conditional law of $\pc$ given $\theta$.
\end{enumerate}

For a window $W \subset \R^d$, our Neyman-Scott clouds are generated by the
standard construction:
    \begin{enumerate}
        \item Expand $W$ by an edge buffer $b$ to produce $W_{exp}$, so that
        clusters whose parents lie outside $W$ are not truncated.
        \item Sample the parent process on $W_{exp}$:
        $n \sim \mathrm{Poisson}(\kappa|W_{exp}|)$, with parents placed uniformly
        on $W_{exp}$.
        \item For each parent $p$ independently, draw an offspring count
        $n_p \sim \mathrm{Poisson}(\mu)$.
        \item For each offspring independently, draw a displacement
        $\epsilon \sim k$ and place the point at $p + \epsilon$.
        \item Discard the parents and retain the offspring lying in $W$.
    \end{enumerate}
The only element that changes from realization to realization is the displacement
kernel $k$. For the Thomas process $k = \mathcal{N}(0, \simpcom^2 I_d)$; for the Mat\'ern
cluster process $k$ is uniform on the ball of radius $R$; for nested Thomas the
parent process at step 2 is itself a Thomas process rather than a homogeneous
Poisson process.

% -----------------------------------------------------------------------------
\subsection{Simulator validation}
\label{ssec:validation}
% -----------------------------------------------------------------------------

% -----------------------------------------------------------------------------
% Commented out for now (2026-08-21): the original moment-matching
% validation study (table + narrative) is hidden per instruction; the
% newer nested-Thomas derivation immediately below (eq:g-nested) is KEPT
% because Table~\ref{tab:classical-comparison}'s footnote $^d$ and three
% Outstanding-experiments bullets cite \eqref{eq:g-nested} directly --
% commenting it out too would leave those as broken "??" references.
% \label{eq:varN} and \label{tab:sim-validation} live inside this hidden
% block; re-enable by deleting the \iffalse/\fi pair below, and note the
% Outstanding-experiments bullet that reproduces
% Table~\ref{tab:sim-validation} for nested Thomas is ALSO commented out
% below for the same reason (it cites this hidden table).
% -----------------------------------------------------------------------------
\iffalse
The validity of the proposed pipeline depends on a diligent cloud-generation
process, which we validate against closed-form second-order theory.

On $W = [0,1]^2$ we compare $3000$ realizations at three parameter triples
spanning the overlap index: tight-cluster ($\nu = 0.07$), moderate
($\nu = 0.22$), and near-Poisson ($\nu = 1.00$).

For the first moment, $\E[N(W)] = \kappa\mu|W|$. For the second, the
canonical identity $\mathrm{Var}[N(W)] = \kappa\mu(1+\mu)|W|$ holds only in the
limit of a window large relative to $\sigma$ and is badly biased in precisely
the near-Poisson regime we wish to stress; the exact finite-window result is
    \begin{equation}
    \label{eq:varN}
    \mathrm{Var}[N(W)] = \kappa\mu + \frac{\kappa\mu^2}{4\pi\sigma^2}J(\sigma)^2,
    \qquad
    J(\sigma) = \int_{-1}^{1}(1-|t|)\,e^{-t^2/4\sigma^2}\,dt,
    \end{equation}
for $W = [0,1]^2$, which collapses to $\kappa\mu(1+\mu)$ as $\sigma \to 0$ and
is the target we compare against, not the canonical identity.

\begin{center}
\captionsetup{type=table}
\small
\begin{tabular}{lccccc}
\toprule
 & & \multicolumn{2}{c}{$\E[N(W)]$} & \multicolumn{2}{c}{$\mathrm{Var}[N(W)]$} \\
\cmidrule(lr){3-4}\cmidrule(lr){5-6}
Case & $\nu$ & theory & empirical & theory & empirical \\
\midrule
A & 0.07 & 200.0 & $200.4 \pm 0.6$ & 982.0 & $972.5 \pm 25.1$ \\
B & 0.22 & 100.0 & $100.2 \pm 0.4$ & 545.2 & $544.8 \pm 14.1$ \\
C & 1.00 & 20.0  & $20.0 \pm 0.1$  & 43.6  & $46.0 \pm 1.2$ \\
\bottomrule
\end{tabular}
\caption{Simulator validation against \eqref{eq:varN}. All deviations are within
$2$ standard errors.}
\label{tab:sim-validation}
\end{center}

We additionally verify that the pair correlation function, Ripley's $K$,
boundary density, and displacement-kernel isotropy all match their
theoretical targets.
\fi

\rev{The same protocol applies to nested Thomas, contrary to what we previously
assumed. Its parent process is a Thomas process with known $g_P$, so $X$ is a Cox
process driven by $\Lambda(u)=\mu_2\sum_{p\in P}k_{\sigma_2}(u-p)$; splitting
$\E[\Lambda(u)\Lambda(v)]$ into diagonal and off-diagonal parts and adding
Gaussian variances under convolution gives $\rho=\kappa\mu_1\mu_2$ and}
    \begin{equation}
    \label{eq:g-nested}
    \rev{g(h) = 1
    + \frac{e^{-\|h\|^2/4\sigma_2^2}}{4\pi\sigma_2^2\,\kappa\mu_1}
    + \frac{e^{-\|h\|^2/4(\sigma_1^2+\sigma_2^2)}}{4\pi(\sigma_1^2+\sigma_2^2)\,\kappa},
    \qquad
    K(t) = \pi t^2
    + \frac{1-e^{-t^2/4\sigma_2^2}}{\kappa\mu_1}
    + \frac{1-e^{-t^2/4(\sigma_1^2+\sigma_2^2)}}{\kappa}.}
    \end{equation}
\rev{This is one clustering term per level, with effective parent intensities
$\kappa\mu_1$ and $\kappa$; as $\sigma_1\to0$ it reduces to a Poisson cluster
process with compound cluster size, $\E[C(C-1)]=\mu_1\mu_2^2+(\mu_1\mu_2)^2$.
% Sentence below cites eq:varN and tab:sim-validation, both inside the
% \iffalse'd-out validation section above -- commented out with it rather
% than left as a broken "??" reference; re-enable together.
% Integrating as in \eqref{eq:varN} gives the finite-window variance, and hence a
% nested-Thomas version of Table~\ref{tab:sim-validation}, which is outstanding.
Two consequences elsewhere: minimum contrast is runnable on nested Thomas, so
its absence from Table~\ref{tab:classical-comparison} is a gap and not an
impossibility; and the derivation needs independent checking before either
claim is relied on.}

% -----------------------------------------------------------------------------
\subsection{Method: ParamNet}
\label{ssec:method}
% -----------------------------------------------------------------------------

Let $\pc$ denote a point cloud, $v$ its vectorized persistence summary, and
$e \in \R^{E + 1}$ the resulting feature vector, where $E$ denotes the encoder's
embedding dimension and the $+1$ is the appended $\log N(\pc)$. ParamNet
follows the pipeline
    \begin{equation}
    \label{eq:pipeline}
    \pc \xrightarrow{\;\text{Vectorization}\;} v
    \xrightarrow{\;\text{Encoder}\;} e
    \xrightarrow{\;\text{FCNN}\;} \hat\theta.
    \end{equation}
We note that the pipeline shape follows the neural-estimator patterns used in
\citet{vihrs2022neural} and \citet{yip2025learning}. Our contribution
lies in the calibration component of the feature-extraction stage, 
the application to a more challenging process (Nested Thomas), an explicit
evaluation of design choices, and comparison of various TDA-based vectorizations.
As mentioned in Section~\ref{sec:background} there is substantial literature
proposing vectorizations of persistence
diagrams -- persistence images, Betti curves, landscapes, kernel methods --
that trade off stability, resolution, and information loss in different ways,
but there has been no formalized discussion on which is preferable in this setting.

ParamNet uses the persistence image vectorization, as it is the best performing (Table~\ref{tab:branches}). We now discuss the key components of the model in turn.

\paragraph{Filtration.} 
Instead of the canonical Rips or \v{C}ech filtrations,
we use the $DTM_k$ filtration (see Example~\ref{ex:dtm}).
Rips is built from inter-point distances alone, making it sensitive to the
sparsest connection between two regions and does not register how densely those
regions are occupied. $DTM_k$ instead works with a point's average distance to its
$k$ nearest neighbours, providing a more direct measure of local density.
Since the parameters of a point process directly affect the local density of the
point cloud, we claim that this choice is natural. [TODO: introduce the idea of using multiple scales, say that the architectural specifics are explained in encoder paragraph, explain that thomas uses k5 and nested thomas uses k5/10/15]

\paragraph{Calibration.}
\label{ssec:calibration}
Fix the homology dimension $h \in \{0,1\}$; we obtain $N_{train}$ persistence
diagrams $\{D^{(h)}_i\}_{i=1}^{N_{train}}$, some of which may have
extreme outliers relative to the full set.
To ensure that in the resulting persistence images each pixel bears the
same meaning, the diagrams must be calibrated to have the same axis ranges. It
is not immediately clear that the naive approach of calibrating to the maximum deaths
and persistences is adequate in this context: if a diagram in the training set has
strong outliers, the remaining diagrams
will be rescaled in a way that places the majority of their points relatively 
close to the origin. [TODO: diagram]

Our solution is to calibrate according to \emph{per-diagram coverage quantiles}.
For each training diagram $D_i^{(h)}$, define the smallest birth
$b_i^{\min}$, largest birth $b_i^{\max}$, and largest persistence value
$p_i^{\max}$ \emph{within a single diagram}. Given a coverage level
$q \in (0,1]$ and a
padding factor $\alpha \geq 1$, we set the calibrated birth and persistence
ranges to be the $q$-quantiles of these per-diagram statistics, taken over the
training set:
\begin{align*}
    \big[\,Q_{1-q}\big(\{b_i^{\min}\}_{i=1}^{N_{train}}\big),\
    \alpha \cdot Q_{q}\big(\{b_i^{\max}\}_{i=1}^{N_{train}}\big)\,\big]
    \qquad \text{and} \qquad
    \big[\,0,\ \alpha \cdot
    Q_{q}\big(\{p_i^{\max}\}_{i=1}^{N_{train}}\big)\,\big],
\end{align*}
where $Q_q(\cdot)$ denotes the empirical $q$-quantile. These two
ranges are shared by every diagram of homology dimension $h$, for every point
cloud in the dataset, so that a
given pixel of the resulting persistence image always refers to the same region
of birth-persistence space.
The calibrated ranges are fit \emph{exclusively} on the training diagrams,
then frozen and applied unchanged to validation and test diagrams;
fitting them on the full, unsplit population of clouds leaks information
about the held-out clouds into the training images.
For $H_0$, birth values are often identically zero under the
chosen filtration, in which case the calibrated birth range
collapses to a point, and we fall back to a
one-dimensional vectorization.
The four calibration parameters -- coverage $q$, padding $\alpha$, kernel width
$\sigma_{\text{pixels}}$, and raster resolution $G$ -- are swept one at a time
around default values. Informed by the calibration sweep, ParamNet takes $q=0.95$, $\alpha=1.05$,
$\sigma_{\text{pixels}}=0.5$ and $G=64$ (values to be re-checked against the run
configuration; $\sigma_{\text{pixels}}$ was changed from an earlier $0.75$).

\paragraph{Encoder and head.} For the Thomas process, we compute the $H_0$ and $H1$ persistence images of the $DTM_5$ diagram. For Nested Thomas, we compute $H_0$ and $H_1$ images for each of the three $DTM_{\{5,10,15\}}$ filtrations and implement a particular architecture to handle the various scales. All images are computed at resolution $64 \times 64$.

For batch size $B$ and $C\in\{1,3\}$ scales, the encoding stage is
\begin{align*}
    (B, C, 64, 64) \xrightarrow{\;\text{Encoder}\;} (B,E \cdot C) \xrightarrow{\;\text{Concatenation}\;} (B,E \cdot C+1)
\end{align*}
In the encoder stage, the image is passed through a small
convolutional stack: three $3\times3$ convolution blocks with channel widths
$(32, 64, 128)$, each followed by a ReLU and $2\times2$ max-pooling with
spatial dropout ($p=0.2$) between blocks. Each block is a CoordConv layer
\citep{liu2018coordconv}: two extra channels holding each pixel's
normalized $x$ and $y$ coordinate are concatenated to the image before the
convolution, giving the otherwise translation-equivariant stack direct
access to absolute position. This matters here because a persistence
image's axes are not arbitrary -- birth encodes local density and
persistence encodes separation -- so a pixel's location, not just its
neighbourhood, is informative. The resulting
feature map is reduced by adaptive average pooling and projected by a linear
layer to the embedding dimension $E$ (we use $E=64$).
The concatenation step appends $\log N(\pc)$, the logged number of points in
the cloud, to that embedding. The
resulting $(B, E + 1)$ tensor is then passed to the FCNN head
$NN_\phi: \R^{E + 1} \to \R^{n_{params}}$, a small feed-forward network with
two hidden layers of width 64 and
32, each followed by a ReLU nonlinearity and dropout ($p=0.1$), then a final
linear layer mapping to the $n_{params}$ outputs.

When three scales are used in Nested Thomas, they share one encoder, which works better than using three independent encoders [TODO: cite evidence]. The three resulting feature vectors in $\R^E$ are fused by simple concatenation.

Every experiment uses the same optimizer and schedule: AdamW ($\text{lr}=10^{-3}$,
weight decay $3\times10^{-4}$), batch size $32$, up to $500$ epochs with
early-stopping patience $50$ on the validation loss.

% -----------------------------------------------------------------------------
\subsection{Evaluation protocol}
\label{ssec:eval}
% -----------------------------------------------------------------------------

Data are split by parameter vector so that no parameter
value appears in both training and evaluation. Beyond the random test split we
generate two further evaluation sets: a grid over the design grid, permitting
error to be reported as a surface over the recalibrated parameter space used in
Section~\ref{ssec:simulator}; and a set drawn
one increment outside each range, characterizing
degradation under extrapolation. The classical baselines are run on the same
clouds and splits rather than on
separate data or numbers taken from the literature.

We report error in the log-normalized
parameter space of Section~\ref{ssec:setup}. For a parameter vector with
$n_{params}$ components and a test set of size $N_{test}$, we define the
per-parameter loss
\begin{align*}
    \mathcal{L}^j = \frac{1}{N_{test}} \sum_{i=1}^{N_{test}}
    (\hat{\tilde\theta}^j_i - \tilde\theta^j_i)^2
    \qquad \text{and} \qquad
    \mathcal{L} = \frac{1}{n_{params}} \sum_j \mathcal{L}^j,
\end{align*}
where $\tilde\theta$ denotes the log-normalized parameter vector and
$\hat{\tilde\theta}$ the model's estimate of it.

We measure the \emph{chance-level} by retraining the default PI model with training
labels randomly permuted (inputs unchanged). Both processes land at
$\mathcal{L}\approx1.0$ (Table~\ref{tab:classical-comparison}), giving an approximate error ceiling against which every other result can be read. 

% -----------------------------------------------------------------------------
\subsection{Estimation performance and baseline comparison}
\label{ssec:results}
% -----------------------------------------------------------------------------

The runs reported below match ParamNet as previously specified on the scale axis: $DTM_5$ for Thomas, fused $DTM_{\{5,10,15\}}$ for nested Thomas. A dagger ($^\dagger$) marks an arm on which at least one seed failed to converge and collapsed instead to the $\mathcal{L}\approx0.9$--$1.0$ chance level; in these cases the reported mean is dominated by the failed seed(s).

We note that training on $\log N$ alone (no
topological summary) collapses to the same
$\mathcal{L}\approx0.9$--$1.0$ chance level as the label-shuffle control
(on both processes), while removing $\log N$ from the
full model costs little on nested-Thomas ($0.192\to0.203$) and is
merely noisy on Thomas rather than clearly harmful (both are the
``Features'' rows of Table~\ref{tab:branches}). Together the two
arms show the persistence images -- not the point count -- are carrying the
signal, with $\log N$ a minor, possibly non-essential addition.

\begin{center}
\captionsetup{type=table}
\small
\begin{tabular}{lcc}
\toprule
Method & Thomas $\mathcal{L}$ & Nested-Thomas $\mathcal{L}$ \\
\midrule
ParamNet$^{\ddagger}$    & $0.125 \pm 0.007$ & $0.192 \pm 0.010^{c}$ \\
vihrs ($L(r)-r$ CNN)     & $0.115 \pm 0.007$ & $0.217 \pm 0.006$ \\
Minimum contrast, $K$   & $0.855 \pm 0.682^{a}$ & \rev{---$^{d}$} \\
Minimum contrast, $g$   & $6.199 \pm 0.339^{b}$ & \rev{---$^{d}$} \\
\midrule
\emph{Chance-level reference (label shuffle)} & \emph{$1.003 \pm 0.019$} & \emph{$1.006 \pm 0.020$} \\
\bottomrule
\end{tabular}
\caption{Estimator comparison on Thomas and Nested Thomas, $n=10$ seeds
unless noted, on the shared test split
of Section~\ref{ssec:eval}. $^{\ddagger}$ParamNet is
process-specific: DTM$_5$ for Thomas, fused
DTM$_{5,10,15}$ for nested Thomas; see Table~\ref{tab:branches}.
$^{c}$$n=3$, the only seeds run at the current calibration
($\sigma_{\text{pixels}}=0.5$); see Section~\ref{ssec:method}.
$^a$ Minimum contrast on
$K$ is unstable across seeds rather than uniformly bad, and $^b$ minimum
contrast on $g$ is consistently worse than chance; both are discussed in
the text. \rev{$^{d}$Not run: the closed form of \eqref{eq:g-nested} makes
this comparison available in principle, so its absence is an outstanding
experiment rather than an impossibility.}}
\label{tab:classical-comparison}
\end{center}

\begin{center}
\captionsetup{type=table}
\small
\textbf{(a) Thomas} ($\kappa$ parent intensity, $\mu$ offspring intensity, $\sigma$ cluster scale)\\[2pt]
\begin{tabular}{lcccc}
\toprule
Parameter & ParamNet & vihrs & Min.\ contrast, $K$ & Min.\ contrast, $g$ \\
\midrule
$\mathcal{L}^\kappa$ & $0.201 \pm 0.010$ & $0.192 \pm 0.013$ & $1.373 \pm 1.087$ & $8.825 \pm 0.485$ \\
$\mathcal{L}^\mu$    & $0.134 \pm 0.007$ & $0.122 \pm 0.007$ & $0.939 \pm 0.758$ & $6.334 \pm 0.370$ \\
$\mathcal{L}^\sigma$ & $0.040 \pm 0.004$ & $0.029 \pm 0.002$ & $0.254 \pm 0.203$ & $3.437 \pm 0.170$ \\
\bottomrule
\end{tabular}

\vspace{10pt}
\textbf{(b) Nested Thomas}\\[2pt]
\begin{tabular}{lcc}
\toprule
Parameter & ParamNet$^{c}$ & vihrs \\
\midrule
parent intensity      & $0.283 \pm 0.023$ & $0.320 \pm 0.010$ \\
meta-offspring        & $0.402 \pm 0.022$ & $0.493 \pm 0.022$ \\
meta-cluster scale    & $0.143 \pm 0.006$ & $0.132 \pm 0.009$ \\
mean offspring        & $0.086 \pm 0.004$ & $0.110 \pm 0.010$ \\
cluster scale         & $0.044 \pm 0.000$ & $0.028 \pm 0.003$ \\
\bottomrule
\end{tabular}
\caption{Per-parameter loss $\mathcal{L}^j$ underlying the aggregate
$\mathcal{L}$ of Table~\ref{tab:classical-comparison} (same methods, same
seeds; ParamNet/nested is $n{=}3$, see that table's $^{c}$). Minimum
contrast has no nested-Thomas columns for the reason given there.}
\label{tab:perparam}
\end{center}

We compare ParamNet against the classical and neural baselines of
Section~\ref{sec:intro}: minimum contrast fit to
Ripley's $K$ and, separately, to the pair correlation function $g$,
and \texttt{vihrs} -- our re-implementation
of the neural estimator of \citet{vihrs2022neural}: a 1-D CNN over the
edge-corrected, centred Ripley $L(r)-r$ curve
plus $\log N(\pc)$, i.e.\ the pipeline of Section~\ref{ssec:method} with the
persistence-image feature replaced by a single classical radial summary.

\texttt{vihrs} beats ParamNet on Thomas: its per-parameter losses
($\mathcal{L}^\kappa=0.192$, $\mathcal{L}^\mu=0.122$,
$\mathcal{L}^\sigma=0.029$) are each individually lower than
ParamNet's own ($0.201$, $0.134$, $0.040$; Table~\ref{tab:perparam}).
So on this single-level process, our topological pipeline does not beat a
hand-built radial summary.

On nested
Thomas, ParamNet's gains an advantage on the density/count targets while
\texttt{vihrs} maintains an edge on both scale parameters. $L(r)-r$ is built to
resolve a characteristic radius directly, while a density-based topological
feature is not, and it is only on the more complex two-level process that this
trade-off favours our pipeline.

Both neural methods dominate minimum contrast on Thomas by roughly
an order of magnitude. The two failure modes are different. On $K$ the estimator achieves losses of $0.365$--$0.399$ for $6/10$ seeds -- close to \texttt{vihrs} -- and $0.903$--$2.056$ on the remaining $4/10$ -- which is the near-Poisson flatness.

On $g$ every seed is worse than chance, indicating
systematic bias rather than scattered non-convergence; the likely cause is our
contrast exponent $c=1$, set by analogy with \citet{diggle2003spatial}'s
$K$-based rule of thumb rather than tuned, where $c=1/4$ is the value validated
for aggregated patterns. We report the number but do not read it as a statement
about contrast estimation on $g$ in general.

Parent intensity $\kappa$ is the hardest target
for every method on both processes by a wide margin and the cluster scale $\sigma$ the
easiest, and this ordering holds
regardless of which estimator is used; so the difficulty ranking looks like a
property of the estimation problem itself rather than
an artifact of any one feature set or architecture.

\paragraph{Branches.} Table~\ref{tab:branches} reports the architectural ablations, each
measured one axis at a time against the multi-scale, $H_0{+}H_1$, CoordConv
configuration this work started from (CoordConv itself -- see
Section~\ref{ssec:method} -- is taken as given rather than ablated). Error grows monotonically with $k$ within a single
scale on both design
distributions: this is consistent with the fact that a finer $k$-nearest neighbourhood
conveys more information about the local-density signal, while a coarser one averages this structure away.
None of
the three alternative vectorizations beats
persistence image on either process or encoder.

\begin{center}
\captionsetup{type=table}
\footnotesize
\begin{tabular}{llcc}
\toprule
Axis & Arm & Thomas $\mathcal{L}$ & Nested-Thomas $\mathcal{L}$ \\
\midrule
Filtration & Rips & ---$^a$ & ---$^a$ \\
           & DTM$_5$ / DTM$_{15}$ & $0.125 \pm 0.007$ / $0.146 \pm 0.006$ & $0.205 \pm 0.003$ / $0.260 \pm 0.017$ \\
           & DTM$_{\{5,10,15\}}$ (fused) & $0.126 \pm 0.008$ & $0.192 \pm 0.010^{c}$ \\
Homology   & $H_0$ only / $H_1$ only & $0.125 \pm 0.006$ / $0.900 \pm 0.019$ & $0.280 \pm 0.250^\dagger$ / $0.932 \pm 0.028$ \\
Encoder    & shared / independent, concat & $0.125 \pm 0.007$ / $0.134 \pm 0.009$ & $0.201 \pm 0.010$ / $0.222 \pm 0.021$ \\
Vectorization & Persistence image & $0.126 \pm 0.008$ & $0.192 \pm 0.010^{c}$ \\
           & Landscape & $0.132 \pm 0.008$ & $0.250 \pm 0.005$ \\
           & Silhouette ($p=1$) & $0.165 \pm 0.013$ & $0.271 \pm 0.012$ \\
           & Betti curve & $0.196 \pm 0.014$ & $0.290 \pm 0.005$ \\
Features   & $\log N$ removed & $0.217 \pm 0.278^\dagger$ & $0.203 \pm 0.007$ \\
           & Topology removed ($\log N$ only) & $0.893 \pm 0.015$ & $0.948 \pm 0.019$ \\
\bottomrule
\end{tabular}
\caption{\new{Branch summary: every axis measured against the multi-scale,
$H_0{+}H_1$, CoordConv starting configuration, at the default $q=0.95$
calibration, $n=10$ seeds per arm unless noted. Vectorization
rows are the native-encoder path. $^a$ Rips fails on $10/10$ seeds before
training starts, for the degenerate-birth reason given in the text.
$^{c}$$n=3$, the only seeds run at the current calibration
($\sigma_{\text{pixels}}=0.5$); see Section~\ref{ssec:method}.}
\rev{The DTM$_5$ / fused rows now support the process-specific scale choice of
Section~\ref{ssec:method} rather than ablating a single model.}}
\label{tab:branches}
\end{center}


% #############################################################################
\section{Discussion, limitations and future work}
\label{sec:future}
% #############################################################################

\new{The following are recorded as notes rather than prose.}

\paragraph{Limitations.}
\begin{itemize}[leftmargin=*,itemsep=2pt]
    \item \rev{\textbf{Two models, not one.} The filtration scale is a
    per-family choice with no rule beyond running both and comparing validation
    loss, which undercuts the process-agnostic framing of
    Section~\ref{sec:intro}. Always fusing would restore a single model at
    $3\times$ the persistence cost and $0.001$ on Thomas, i.e.\ inside seed
    noise; attention over scales would be the principled fix.}
    \item \textbf{Negative results.} Three alternative vectorizations were
    implemented and none improved on persistence images. This is evidence about
    these summaries under this architecture and
    this design distribution, not a general claim about vectorizations; the
    same qualification covers $H_1$, which fails to earn its place \emph{here}.
    \item \rev{\textbf{No advantage on Thomas.} \texttt{vihrs} wins on every
    parameter of the single-level process, and on both scale parameters of the
    nested one. Claims about topological features should be stated at that
    resolution.}
    \item \textbf{Calibration inherited across summaries.} The per-diagram
    coverage rule was inherited unchanged by the 1-D summaries for
    comparability, which may not be optimal for each individually.
    \item \textbf{Single realization; homogeneity.} The estimator sees one
    realization, and everything here assumes
    homogeneous, stationary, isotropic processes.
    \item \textbf{The design distribution is a prior.} The design distribution
    acts as a prior, so the trained model approximates a posterior summary
    under it, and its ranges must therefore cover the intended application
    domain -- the single most important limitation of any amortized
    estimator. \rev{\citet{vihrs2022neural} gives the practical check: verify
    that the observed point count and summary curve fall inside the envelope of
    the training simulations before trusting an estimate.}
    \item \textbf{Computational cost.} Persistence computation scales badly in
    $n$; this is why the point budget of Section~\ref{ssec:simulator} is capped at
    800. \new{The amortized inference cost per cloud is currently unmeasured.}
    \rev{\citet{vihrs2022neural} reports the comparison to make: seconds to fit
    all 5000 test patterns with the network, against $\sim$1.2 s \emph{each} for
    minimum contrast.}
\end{itemize}

\paragraph{Outstanding experiments.} \new{The following results are promised
elsewhere in this report but not yet in hand; they are the nearest-term
deliverable.}
\begin{itemize}[leftmargin=*,itemsep=2pt]
% Two bullets below removed (2026-08-21): both were "outstanding run"
% references -- the nested-Thomas fused fill-in (n=3->10) and a stale
% headline-config re-statement (still said "H0 only" / "no coordinate
% channels", both superseded -- see Sec ssec:method and the CoordConv
% paragraph there). Current best data is used directly in the tables now
% instead of being flagged as pending.
\iffalse
    \item \rev{Complete the nested-Thomas fused run at $n=10$ seeds; the
    two-model split rests entirely on a provisional $n{=}3$ cell.}
    \item Train and evaluate the headline configuration exactly as
    Section~\ref{ssec:method} specifies it -- persistence image, $H_0$ only,
    \rev{$DTM_5$ (Thomas) / fused (nested)}, shared encoder,
    concatenation with $\log N$, no coordinate channels, $q=0.95$ -- on both design
    distributions at $n=10$ seeds, so that the reported $\mathcal{L}$ is the model
    the report describes.
\fi
% Bullet below commented out with the irreducible-error-floor paragraph in
% ssec:setup (2026-08-21) -- report stands without it for now.
\iffalse
    \item \textbf{The error floor.} $\mathcal{L}_{\text{floor}}$ does not yet
    have a value to report: the estimation script was run
    on a dedicated $50\,\theta \times 200$-realization Thomas dataset and every
    shard completed without error, but every per-cloud loss reduced to
    \texttt{NaN} rather than a number. One concrete lead: the shard JSONs'
    recorded \texttt{label\_names} include names that do not belong to
    Thomas's own
    $(\kappa,\mu,\sigma)$ target set, suggesting the per-target loss array is
    being indexed against the wrong label list. Fix the
    label-indexing bug and re-run for Thomas\rev{; nested Thomas needs no
    separate treatment, since neither route in Section~\ref{ssec:setup} requires
    a closed form and \eqref{eq:g-nested} now supplies one anyway.}
\fi
    \item \rev{Run minimum contrast on nested Thomas against
    \eqref{eq:g-nested}, and re-run it on $g$ with the exponent $c$ swept.}
% Bullet below commented out (2026-08-21): cites Table~\ref{tab:sim-validation},
% which lives inside the now-hidden Simulator-validation section.
\iffalse
    \item \rev{Validate the nested-Thomas simulator against
    \eqref{eq:g-nested}, reproducing Table~\ref{tab:sim-validation}.}
\fi
    \item Three evaluation sets defined in Section~\ref{ssec:eval} are not
    yet reported: error as a surface
    over $(\log\nu, \log\lambda)$ on the grid set -- expect degradation as
    $\nu \to 1$, and if the model
    degrades more gracefully than minimum contrast there, that IS the result;
    degradation on the out-of-range set; and the Mat\'ern cluster runs for the
    cross-family comparison -- whether one
    model trained across families beats per-family models is an interesting question
    in itself, and the common reparameterization of
    Section~\ref{ssec:simulator} is what makes it askable.
    \rev{Mat\'ern is also the family \citet{waagepetersen2009twostep} exclude on
    smoothness grounds, so it is the natural place to show a simulation-trained
    estimator has no such restriction.}
    \item \new{Run the Palm likelihood baseline \citep{tanaka2008parameter} to
    completion; it is implemented but has not yet been run.}
\end{itemize}

\paragraph{Planned directions.}
\begin{itemize}[leftmargin=*,itemsep=2pt]
    \item \textbf{Near term.} The encoder question: attention over scales, and
    set-based encoders operating on diagrams directly. \rev{The former would also
    dissolve the two-model limitation.}
    \item \textbf{Medium term.} Inhomogeneous / non-stationary processes, where
    classical closed forms are least available and the topological approach has
    the most to gain -- this is arguably where the thesis-scale contribution
    lies. \rev{The comparator is the two-step procedure of
    \citet{waagepetersen2009twostep}, which needs the covariates observed and the
    intensity correctly specified; a simulation-trained estimator needs neither.}
    \item \textbf{Medium term.} Real data. Galaxy catalogues if the running
    example of Section~\ref{sec:intro} is cosmological.
    \rev{The Barro Colorado Island tree data of \citet{waagepetersen2009twostep}
    is a good second target: small, public, and already fitted with a Thomas
    process.}
    \item \textbf{Longer term.} Uncertainty quantification -- the model
    currently emits a point estimate with no interval, which limits its use as
    a drop-in replacement for likelihood-based methods.
    \rev{This is the same machinery as route (ii) of Section~\ref{ssec:setup}.}
\end{itemize}


% #############################################################################
% REFERENCES
% #############################################################################
{\footnotesize
\setlength{\bibsep}{1pt plus 0.3ex}
\begin{thebibliography}{9}
\setlength{\itemsep}{0pt}
\setlength{\parskip}{0pt}
\bibitem[Adams et al.(2017)]{adams2017persistence}
H.~Adams et al.
\newblock Persistence images: a stable vector representation of persistent
homology.
\newblock \emph{Journal of Machine Learning Research}, 18(8):1--35, 2017.

\bibitem[Anai et al.(2020)]{anai2020dtm}
H.~Anai, F.~Chazal, M.~Glisse, Y.~Ike, H.~Inakoshi, R.~Tinarrage, and
Y.~Umeda.
\newblock DTM-based filtrations.
\newblock In \emph{Topological Data Analysis}, Abel Symposia 15, pages 33--66.
Springer, 2020.

\bibitem[Biagetti et al.(2021)]{biagetti2021persistence}
M.~Biagetti, A.~Cole, and G.~Shiu.
\newblock The persistence of large scale structures.
\newblock \emph{Journal of Cosmology and Astroparticle Physics}, 2021(04):061,
2021.

\bibitem[Bubenik(2015)]{bubenik2015landscapes}
P.~Bubenik.
\newblock Statistical topological data analysis using persistence landscapes.
\newblock \emph{Journal of Machine Learning Research}, 16:77--102, 2015.

\bibitem[Chazal et al.(2014)]{chazal2014silhouette}
F.~Chazal, B.~T. Fasy, F.~Lecci, A.~Rinaldo, and L.~Wasserman.
\newblock Stochastic convergence of persistence landscapes and silhouettes.
\newblock In \emph{Proc.\ 30th Annual Symposium on Computational Geometry},
pages 474--483, 2014.

\bibitem[Diggle(2003)]{diggle2003spatial}
P.~J. Diggle.
\newblock \emph{Statistical Analysis of Spatial Point Patterns}.
\newblock Arnold, London, 2nd edition, 2003.

\bibitem[Liu et al.(2018)]{liu2018coordconv}
R.~Liu et al.
\newblock An intriguing failing of convolutional neural networks and the
CoordConv solution.
\newblock In \emph{Advances in Neural Information Processing Systems}, 2018.

\bibitem[M\o ller and Waagepetersen(2004)]{moller2004statistical}
J.~M\o ller and R.~P. Waagepetersen.
\newblock \emph{Statistical Inference and Simulation for Spatial Point
Processes}.
\newblock Chapman \& Hall/CRC, 2004.

\bibitem[Reininghaus et al.(2015)]{reininghaus2015kernel}
J.~Reininghaus, S.~Huber, U.~Bauer, and R.~Kwitt.
\newblock A stable multi-scale kernel for topological machine learning.
\newblock In \emph{IEEE Conference on Computer Vision and Pattern Recognition},
pages 4741--4748, 2015.

\bibitem[Tanaka et al.(2008)]{tanaka2008parameter}
U.~Tanaka, Y.~Ogata, and D.~Stoyan.
\newblock Parameter estimation and model selection for Neyman-Scott point
processes.
\newblock \emph{Biometrical Journal}, 50(1):43--57, 2008.

\bibitem[Umeda(2017)]{umeda2017time}
Y.~Umeda.
\newblock Time series classification via topological data analysis.
\newblock \emph{Transactions of the Japanese Society for Artificial
Intelligence}, 32(3):1--12, 2017.

\bibitem[Vihrs(2022)]{vihrs2022neural}
N.~Vihrs.
\newblock Using neural networks to estimate parameters in spatial point process
models.
\newblock arXiv preprint arXiv:2109.15056, 2022.

\bibitem[Waagepetersen and Guan(2009)]{waagepetersen2009twostep}
R.~Waagepetersen and Y.~Guan.
\newblock Two-step estimation for inhomogeneous spatial point processes.
\newblock \emph{Journal of the Royal Statistical Society: Series B},
71(3):685--702, 2009.

\bibitem[Yip et al.(2025)]{yip2025learning}
J.~H.~T. Yip, A.~Rouhiainen, and G.~Shiu.
\newblock Learning from topology: cosmological parameter estimation from the
large-scale structure.
\newblock \emph{Machine Learning: Science and Technology}, 6:025063, 2025.
\end{thebibliography}
}

\end{document}